\chapter{Серия \texttt{8a--9a}: транспортная стена, similarity и достройка путей}

\section{Почему после \texttt{7x} уже нельзя было просто искать еще один denominator}

К началу \texttt{8a} две надежды уже были достаточно честно проверены:

\begin{enumerate}[leftmargin=2em]
\item post-accumulation denominators сами по себе не отделяют методы на реальном
эмбрионе;
\item overlap cue-fields тоже не создает specificity там, где transport до цели
не доходит.
\end{enumerate}

Следовательно, следующий вопрос должен был звучать иначе:

\begin{quote}
может быть, проблема не в том, как пересчитать уже полученное поле, а в том, что
сигнал \emph{вообще идет не туда} или не доходит до нужного manifold?
\end{quote}

\section{\texttt{exp8a}: bypass и repair как две попытки пробить transport wall}

\subsection{Замысел}

\texttt{8a} сознательно отвернулся от еще одной ``кухни знаменателей'' и взялся
за два других рычага:

\begin{enumerate}[leftmargin=2em]
\item \textbf{bypass:} дать similarity прямой доступ в итоговый score, а не
держать ее только как gate внутри diffusion;
\item \textbf{repair:} ослабить смазывающие типы ребер, в первую очередь
\texttt{RelatedTo}, чтобы transport меньше стекал в catch-all zones.
\end{enumerate}

Появилась и еще одна методологическая новация: вместо проверки одного
канонического target node вводится \textbf{target manifold}. Это очень важный
сдвиг. С этого места серия перестает считать провалом любой случай, когда граф
выдал не exact label, а корректный соседний graph-native ответ.

\subsection{Методы}

Испытывались следующие варианты:

\[
M0(v)=\frac{transport(v)}{\PR(v)^\eta},
\]
\[
M_{\mathrm{gate}}(v)=transport_{\mathrm{gated}}(v),
\]
\[
M1(v)=\frac{transport(v)+\lambda\,\simf(v\mid q)}{\PR(v)^\eta},
\]
\[
M2(v)=\frac{transport_{\mathrm{routed}}(v)+\lambda\,\simf(v\mid q)}{\PR(v)^\eta},
\]
\[
M3(v)=\frac{M2(v)}{\text{genericness}(v)^\mu}.
\]

В качестве control сохранялся и overlap-вариант из \texttt{7c}.

\subsection{Почему сравнение было честным}

Особенно важно, что \(M0 \to M1\) меняет только одну вещь: additive similarity.
Это позволяет отделить старый semantic gate как ``legacy control'' от нового
вопроса: помогает ли similarity, если дать ей прямой канал в score.

\subsection{Что именно проверялось}

Primary suite была достаточно строгой:

\begin{itemize}[leftmargin=2em]
\item \texttt{coffee + stimulant} против \texttt{beverage}-family;
\item \texttt{check in + document} против procedural and generic distractors;
\item \texttt{laptop + battery} против \texttt{office}/\texttt{power};
\item \texttt{wake up + tired} против \texttt{morning}/\texttt{bed}.
\end{itemize}

Параллельно были control queries, где broad topic не должен был быть уничтожен
over-correction.

\subsection{Главный результат}

В своей окончательной, уже подправленной версии \texttt{8a} дал очень чистую
картину:

\begin{center}
\begin{tabular}{lcccc}
\toprule
Метод & primary manifold & strict & mean rank & control \\
\midrule
\texttt{M0} & 2/4 & 1/4 & 232.0 & 2/2 \\
\texttt{M\_gate} & 2/4 & 1/4 & 217.2 & 2/2 \\
\texttt{M1} & 2/4 & 1/4 & 63.5 & 2/2 \\
\texttt{M2} & 2/4 & 1/4 & 63.5 & 2/2 \\
\texttt{M3} & 2/4 & 1/4 & 135.8 & 2/2 \\
\texttt{CTRL\_overlap} & 1/4 & 0/4 & 5021.2 & 2/2 \\
\bottomrule
\end{tabular}
\end{center}

\subsection{Что из этого следует}

Здесь одновременно есть сильный отрицательный и умеренный положительный
результат.

\subsubsection*{Сильный отрицательный результат}

Overlap control теперь уже не просто ``не победил'', а стал явно хуже. Значит,
ветка ``многократное схождение cue-fields само по себе родит правильный
ответ'' может считаться практически исчерпанной на данном субстрате.

\subsubsection*{Умеренный положительный результат}

Additive similarity действительно улучшает rank structure. Это видно по падению
mean rank с 232.0 до 63.5 у \(M1\). Но она не увеличивает число дискретных
побед.

Иными словами, bypass способен подтягивать правильный manifold из глубины хвоста,
но пока не превращается в механизм надежного переворачивания hardest pairwise
cases.

\subsubsection*{Еще один отрицательный результат}

Routing через \texttt{RelatedTo} в текущем виде почти ничего не прибавил:
\(M2\) практически совпал с \(M1\). Genericness penalty в \(M3\) даже ухудшила
часть показателей.

\subsection{Качественное чтение отдельных запросов}

Именно top-k inspection делает выводы \texttt{8a} убедительными.

\begin{itemize}[leftmargin=2em]
\item В coffee-case \texttt{caffeine} действительно становится ближе, но над ним
остаются \texttt{served hot}, \texttt{coffee shop}, \texttt{made from coffee
beans}.
\item В passport-case система с additive similarity не находит \texttt{passport},
а проваливается в морфологические и procedural neighbours вроде
\texttt{documentor}, \texttt{photodocument}, \texttt{checking in}.
\item В charger-case signal по-прежнему вязнет в \texttt{power}-family и
соседнем structural noise.
\item Morning-case выглядит лучше других: \texttt{drink coffee} действительно
входит в сильный manifold, но не кристаллизуется как бесспорный победитель.
\end{itemize}

\begin{takeaway}
\texttt{8a} не нашел ``победную формулу'', но очень качественно расчистил поле.
После него уже можно уверенно говорить: проблема не решается ни overlap, ни еще
одним знаменателем; additive similarity полезна, но работает как мягкий rescue,
а не как полноценный retrieval engine.
\end{takeaway}

\section{\texttt{exp8b}: de-coning similarity как проектная развилка}

\subsection{Статус}

На момент этого свода \texttt{8b} существует как тщательно продуманная ветка, но
не как исполненный ноутбук. Это нужно проговорить прямо.

\subsection{Почему эта ветка возникла}

Как только similarity из \texttt{8a} становится first-class term, встает вопрос:
а хороша ли вообще сама similarity. Если embedding-space анизотропно и живет в
узком конусе, то cosine margins изначально сжаты, а generic nodes становятся
embedding-hubs.

Тогда часть дефекта, который мы приписываем transport, на самом деле может
оказаться следствием грязной геометрии \texttt{NumberBatch}.

\subsection{Что именно предполагалось проверить}

\texttt{8b} должен был сравнить:

\begin{itemize}[leftmargin=2em]
\item raw cosine;
\item mean-centering;
\item all-but-the-top;
\item CSLS;
\item per-query z-score.
\end{itemize}

Причем сначала на intrinsic критериях (pairwise cosine distribution, hubness,
margin widening), а затем как plug-in replacement для \(\simf(v\mid q)\) в
лучшем score из \texttt{8a}.

\subsection{Почему даже неисполненная ветка все равно важна}

Потому что она задает очень важное ограничение для \texttt{9a}: если
semantic-proximity edges строить на сыром косинусе в узком конусе, можно просто
переписать embedding-hubness в topology и тем самым ухудшить граф.

\begin{openbox}
Если бы \texttt{8b} дал сильный положительный результат, это сделало бы
de-coned similarity стандартным строительным материалом и для additive bypass, и
для semantic-edge construction. Если бы не дал, это сняло бы важный конфаундер и
позволило бы еще увереннее винить transport rather than cosine. Именно поэтому
ветка проектно значима, даже еще не будучи прогнанной.
\end{openbox}

\section{\texttt{exp9a}: dual-geometry transport и semantic-proximity edges}

\subsection{Большой вопрос}

Если \texttt{8a} показал, что additive similarity лишь подтягивает правильные
узлы, но не решает проблему содержательно, то возникает более радикальная идея:
может быть, в графе просто \emph{нет нужных путей}, и их надо аккуратно
достроить.

Отсюда и dual-geometry picture:

\begin{itemize}[leftmargin=2em]
\item graph geometry: typed relational edges;
\item embedding geometry: proximity in vector space.
\end{itemize}

Новая идея состояла в том, чтобы добавить minority channel semantic edges:
mutual-kNN на embeddings, но с mutual filter, degree caps и throttled
conductance.

\subsection{Постановка}

Испытывались семейства:

\begin{itemize}[leftmargin=2em]
\item \(G0b\): no edges, no bypass;
\item \(G1,G2,G3\): только semantic edges с растущим \(w_{\mathrm{sem}}\);
\item \(G0a\): realistic baseline из \texttt{8a} с additive similarity;
\item \(G1r,G2r,G3r\): edges + bypass.
\end{itemize}

Параллельно измерялись два danger checks:

\begin{itemize}[leftmargin=2em]
\item hub reflood: generic distractors не должны снова полезть наверх;
\item RAG collapse: ranking не должен выродиться в pure embedding-kNN.
\end{itemize}

\subsection{Главная сводка}

Результат по основным retrieval-запросам оказался таким:

\begin{center}
\begin{tabular}{lccc}
\toprule
Метод & \((w_{\mathrm{sem}}, \lambda)\) & primary pairwise & median rank \\
\midrule
\texttt{G0b} & (0.0, 0.0) & 1/3 & 249 \\
\texttt{G1} & (0.1, 0.0) & 1/3 & 127 \\
\texttt{G2} & (0.3, 0.0) & 1/3 & 90 \\
\texttt{G3} & (1.0, 0.0) & 1/3 & 43 \\
\texttt{G0a} & (0.0, 1.0) & 1/3 & 111 \\
\texttt{G1r} & (0.1, 1.0) & 1/3 & 98 \\
\texttt{G2r} & (0.3, 1.0) & 1/3 & 59 \\
\texttt{G3r} & (1.0, 1.0) & 0/3 & 42 \\
\bottomrule
\end{tabular}
\end{center}

\subsection{Как это читать}

Здесь снова есть важное разделение между rank-improvement и task-solution.

\begin{itemize}[leftmargin=2em]
\item Да, semantic edges заметно подтягивают target manifolds вверх по рангу.
\item Нет, они не увеличивают число дискретных побед над strongest distractor.
\item При слишком сильном semantic-channel combined regime даже портится
(см. \(G3r\)).
\end{itemize}

\subsection{Сильная сторона \texttt{9a}}

Особенно интересным оказался bridge probe на alias/synonym pairs:

\begin{itemize}[leftmargin=2em]
\item \texttt{cat -> kitty}: rank 1161 \(\to\) 285 \(\to\) 278;
\item \texttt{couch -> sofa}: 1488 \(\to\) 37 \(\to\) 22.
\end{itemize}

То есть там, где проблема действительно состоит в отсутствии короткого моста
между очень близкими представлениями, semantic channel делает реальную работу.

\subsection{Почему этого оказалось недостаточно}

В сложных retrieval cases, особенно passport-case, semantic edges начинали
соединять не ``правильный ответный смысл'', а лексико-морфологическое облако:
\texttt{checking in}, \texttt{check in counter}, \texttt{documentor},
\texttt{photodocument}, \texttt{redocument} и т.\,п.

Это очень сильный диагностический результат. Он показывает, что blind
embedding-proximity repair шьет graph gaps, но шьет их не по той семантике,
которая нужна ассоциативной памяти.

\subsection{Danger checks}

Здесь ноутбук тоже оказался очень полезен.

\begin{itemize}[leftmargin=2em]
\item Hub reflood не взорвался на умеренных \(w_{\mathrm{sem}}\): median rank
generic distractors даже немного рос.
\item Но RAG-collapse стал заметен: overlap top-20 с pure-sim ranking вырос
примерно от 0.10 до 0.52--0.53.
\end{itemize}

Это значит, что сильный semantic channel начинает подменять граф векторным
retrieval. Память при этом не столько улучшается, сколько тайком становится
другим механизмом.

\begin{takeaway}
\texttt{9a} не дал структурного ``спасения'' линии, но очень хорошо локализовал
предел blind semantic-edge repair. Он показал, что у эмбриона действительно не
хватает мостов, однако строить их как глобальные mutual-kNN links по сырой
NumberBatch-геометрии недостаточно: такие мосты тянут не только к смысловым
соседям, но и к морфологическим и лексическим паразитным облакам.
\end{takeaway}

\section{Что серия \texttt{8a--9a} установила}

\subsection{Отрицательные результаты}

\begin{enumerate}[leftmargin=2em]
\item Overlap cue-fields не является недостающим ингредиентом retrieval.
\item Ни relation-type routing в текущем виде, ни genericness penalty не дали
решающего эффекта.
\item Blind semantic-proximity edges не превращают эмбрион в хорошую память.
\end{enumerate}

\subsection{Положительные и диагностические результаты}

\begin{enumerate}[leftmargin=2em]
\item Additive similarity полезна как мягкий bypass: она улучшает rank structure
и поднимает цель из глубокого хвоста.
\item Semantic edges годятся как средство alias/synonym bridging.
\item Но в содержательных commonsense retrieval cases главная проблема остается
глубже: signal либо идет по неправильным typed routes, либо graph-native answer
organizes itself not as the named target node.
\end{enumerate}

\subsection{Почти окончательный вывод этой линии}

К концу \texttt{9a} становится правдоподобной уже достаточно смелая мысль:

\begin{quote}
в текущем виде диффузионная retrieval-парадигма, даже с priors, leak laws,
overlap and semantic repair, может быть уперлась в свой естественный потолок на
данном субстрате.
\end{quote}

Это не означает, что диффузия бесполезна. Это означает, что от нее, возможно,
требуют того, что ближе к path-selective reasoning или к архитектуре
консолидации графа, чем к простому spreading process.
